\section{光的干涉}
\label{sec:光的干涉}
在\xref{sec:光的电磁理论}中，我们在波动光学重新认识了光的基本形态，并引入了光强和光程这两个非常重要的概念。而本节我们就来思考这样一个问题，如果两束光在某处叠加将会导致什么？类比机械波的知识，我们知道，两列同频率且具有固定相位差的波将会发生干涉现象，而这种干涉现象是一切波动所具有的共同特性，无所谓是机械波还是电磁波，因此光也会发生干涉现象。

\subsection{光的相干性}
\label{subsec:光的相干性}
光并不总是会发生干涉现象的，而是要满足以下的\uwave{相干条件}
\begin{enumerate}
    \item 光的频率相同
    \item 光存在平行的振动分量
    \item 光在叠加时的相位差恒定
\end{enumerate}
我们将满足相干条件的光波称为\uwave{相干光波}，在上述三个条件中，第二条是比较费解的，这是因为我们知道，取决于光的传播方向，光矢量的振动方向是不一样的，试想若两束光的光矢量分别在$x$轴和$y$轴方向振动，即便它们频率相同且相位差恒定，那显然也无法产生干涉现象。

我们现在来思考，光的干涉将导致什么现象？我们可能觉得不会有什么现象，因为干涉不过是简谐振动的叠加，而简谐振动的叠加仍然是简谐振动。但问题在于，这种叠加导致了振动在空间上的分布，也就是波动的形态发生了变化，有些点的振动加强，有些点的振动减弱，而我们知道我们能看到的光的明暗实质是光强，而光强又与光振幅的平方成正比，因此我们就可以预期，光发生干涉时，有些地方会更亮，有些地方会更暗，\empx{光强在干涉时并不是简单的叠加}。

在这一小节，我们姑且先考虑光的干涉在某一点的光强，假设空间中有两束同频率、相位差恒定、光振动方向一致的光波，传播到空间中某点$P$处，记其分别引发的光振动为$E_A,E_B$，则
\gdef\prefixeq{光的相干性}
\begin{Gather}[6pt]
    E_\text{A}=E_{\text{A}0}\cos(\omega t+\phi_{\text{A}})\xlabelpeq{1}\\
    E_\text{B}=E_{\text{B}0}\cos(\omega t+\phi_{\text{A}})\xlabelpeq{2}
\end{Gather}

根据光矢量的叠加原理，在任意时刻都有
\begin{Equation}&[3]
    E=E_1+E_2
\end{Equation}
其中$E$为点$P$处光的合振动，而两个简谐振动的合成仍然为简谐振动
\begin{Equation}&[4]
    E=E_0\cos(\omega t+\phi)
\end{Equation}
其中上式光的合振动的振幅为\footnote{这里合振动的初相位$\phi$也是可以求出的，其为$\mal{\phi=\arctan\frac{E_{\text{A}0}\sin\phi_{\text{A}}-E_{\text{B}0}\sin\phi_{\text{B}}}{E_{\text{A}0}\cos\phi_{\text{A}}-E_{\text{B}0}\cos\phi_{\text{B}}}}$，但此处用不到该结果。}
\begin{Equation}&[5]
    E_0=\sqrt{E_{\text{A}0}^2+E_{\text{B}0}^2+2E_{\text{A}0}E_{\text{B}0}\cos(\phi_{\text{B}}-\phi_{\text{A}})}
\end{Equation}
而在这种情况下，根据\fancyref{fml:光的强度}
\begin{Equation}&[6]
    I=\frac{n}{\mu_0c}\frac{1}{\tau}\Int[0][\tau]{E^2\dd{t}}=\frac{n}{\mu_0c}\frac{1}{\tau}\Int[0][\tau]{E_0^2\cos(\omega t+\phi)\dd{t}}
\end{Equation}
这就得到，其中$\delt{\phi}=\phi_\text{B}-\phi_\text{A}$为相位差
\begin{Equation}&[7]
    I=\frac{n}{2\mu_0c}\qty[E_{\text{A}0}^2+E_{\text{B}0}^2+2E_{\text{A}0}E_{\text{B}0}\cos\delt{\phi}]
\end{Equation}
我们考虑到$E_{\text{A}0}^2(n/2\mu_0 c)=I_{\text{A}}$以及$E_{\text{B}0}^2(n/2\mu_0 c)=I_{\text{B}}$
\begin{OldBoxFormula}[简谐光的干涉光强]
    简谐相干光在发生干涉时的光强为，其中$\delt{\phi}$为相位差
    \begin{Equation}&[8]
        I=I_{\text{A}}+I_{\text{B}}+2\sqrt{I_\text{A}}\sqrt{I_\text{B}}\cos{\delt{\phi}}
    \end{Equation}
    这就是说，干涉时的光强是两束光的光强之和，再加上$2\sqrt{I_\text{A}}\sqrt{I_\text{B}}\cos{\delt{\phi}}$，称为\uwave{干涉项}。

    如果定义\uwave{光程差} $\delta=\Delta_\text{B}-\Delta_{A}$，那么考虑到$\delt{\phi}=k\delta=(2\pi/\lambda)\delta$
    \begin{Equation}
        I=I_{\text{A}}+I_{\text{B}}+2\sqrt{I_\text{A}}\sqrt{I_\text{B}}\cos(\frac{2\pi}{\lambda}\delta)
    \end{Equation}
    这样就将光强的干涉项与相位差的关系转化为与光程差的关系。
\end{OldBoxFormula}

由\xref{fml:简谐光的干涉光强}中我们不难看出相位差和光程差与\uwave{干涉加强处}和\uwave{干涉减弱处}的关系
\begin{OldTable}[干涉光强与相位差和光程差的关系]
<状态&\mc{2}{相位差}&\mc{2}{光程差}&取值\\>
干涉加强&$\pi$的偶数倍&$\delt{\phi}=\pm \pi(2m)$&半波长的偶数倍&$\delta=\pm(\lambda/2)(2m)$&$m\in\N$\\
干涉减弱&$\pi$的奇数倍&$\delt{\phi}=\pm \pi(2m-1)$&半波长的奇数倍&$\delta=\pm(\lambda/2)(2m-1)$&$m\in\N^{*}$\\
\end{OldTable}

特别的，如果参与干涉的两束光的光强相同，即$I_\text{A}=I_\text{B}=I_0$，那么\xref{fml:简谐光的干涉光强}有
\begin{Equation}
    I=2I_0+2I_0\cos(\frac{2\pi}{\lambda}\delta)
\end{Equation}
这时干涉加强处将有$I=4I_0$即达到单个光波光强的四倍，而干涉减弱处则有$I=0$。

\subsection{光源的发光机制}
\label{subsec:光源的发光机制}
在\xref{subsec:光的相干性}中我们讨论了光的干涉现象，但这总让人感觉有些奇怪，毕竟当我们打开桌上的台灯时，我们的桌子将被均匀的照亮，我们从未在该过程中看到过任何干涉现象，尽管从理论上来说灯管上两个不同位置的光源似乎是可以发生干涉的，这就与光源的发光机制有关了。

在化学中，我们曾经提到过，电子由激发态跃迁至基态时将以电磁波的形式放出能量，这也就是物质燃烧时火焰的光亮来源。而实际上，现代使用的各类照明设施，无论是节能灯（荧光灯和电子镇流器）还是LED灯（发光二极管），尽管用于激发电子的方式不同，但是就其发光的实质而言仍然是电子跃迁释放的可见光频段的电磁波\footnote{总之通过RLC电路的电磁振荡是无法产生可见光频段的电磁波的，因为可见光的频率太高了。}。而我们知道，尽管电子跃迁是大量且持续的，使得宏观上能产生稳定的发光现象，但是单个电子跃迁的持续发光时间都是非常短的，因此光波的初相位是不断变化的（每换一个电子就换一个初相位），而位于光源两个不同位置的光波往往是由两个不同的原子发射的，由于初相位均在不断变化，因此两者的光波在空间相遇时就不具有稳定的相位差，这样就无法形成干涉现象。在数学上就是$\cos{\delt{\phi}}$并不是一个无关时间的常数，不能从\xrefpeq[光的相干性]{6}中作为常数提出，而由于相位差$\delt{\phi}$的变化是完全随机的，因此$\cos{\delt{\phi}}$通过积分求出的均值为零，这样干涉项就为零了，此时的总光强即为
\begin{Equation}
    I=I_1+I_2
\end{Equation}
即光在非相干叠加时的光强就是两者光强的和，这个结果也很好理解，原先在干涉叠加时振动增强和振动减弱是稳定的，这就产生了在时间平均上不会消失的干涉项，而现在这种增强和减弱是不稳定的，能量在空间上就平均分配了，总的光强也就自然是两束光的光强之和了。

此外需要注意的是，普通的光源发出的光波是包含不同波长的复色光，在实践中我们可以通过各种方式将光的频率波长在一个小的范围内，形成适合干涉的准单色光。另外一种更为先进的作法是使用激光，它的原理是通过外加辐射控制电子跃迁，从而产生单色性很好的光源。

总而言之，光的干涉的条件相较于机械波或低频的无线电波的干涉要苛刻的多，光的单色性还姑且较容易解决，关键是在于是如何得到了两列稳定的相位差的光波，为此科学家设计了各类干涉装置，其共同思想是，先将同一处光源发出的光波分为两束，然后使两束光经过不同的光路再在空间中重新相遇产生干涉。取决于具体的分光方式，有两种不同的干涉方式
\begin{itemize}
    \item 第一种方法称为\uwave{分波面干涉法}，以\uwave{双缝干涉}为代表。
    \item 第二种方法称为\uwave{分振幅干涉法}，以\uwave{薄膜干涉}为代表。
\end{itemize}
在接下来的两节中，我们将分别介绍这两种干涉方法。
